\section{Introduction}

Water-quality monitoring in remote Amazonian areas is constrained by sparse infrastructure, high field logistics cost, and intermittent network connectivity. In this context, continuous dependence on cloud inference is operationally fragile and energetically expensive, especially for battery-powered deployments and long-duration campaigns. At the same time, water-quality variables such as pH, dissolved oxygen (OD), and turbidity are highly relevant for environmental screening and public-health risk mitigation.

Conventional machine-learning pipelines can provide strong predictive performance, but many are designed for server-centric execution and do not directly address energy constraints at the edge. Neuromorphic models based on spiking neural networks (SNNs) offer an alternative computational paradigm with event-driven processing and potential compatibility with low-power embedded systems. However, practical evidence in real environmental sensing workflows remains limited, particularly when combining multimodal prediction and platform-level energy benchmarking.

This paper proposes and evaluates an offline multimodal SNN pipeline for water-quality prediction in an Amazonian field scenario, including end-to-end deployment and cross-platform energy analysis. The study integrates data processing, model training, geospatial interpretation, and embedded inference on heterogeneous hardware (ESP32, Raspberry Pi, and Android devices).

The main contributions of this work are an offline multimodal SNN approach for joint prediction of pH, OD, and turbidity from georeferenced physicochemical inputs, a reproducible energy-measurement protocol based on INA219 incremental idle/load estimation across embedded and mobile platforms, a comparative latency--energy analysis that exposes practical trade-offs between microcontroller, SBC, and smartphone execution, and an applied environmental case study in the Amazon region with spatial visualization and operational pollution-index mapping.
